% chapters/abstract.tex
\chapter*{Abstract}
\addcontentsline{toc}{chapter}{Abstract}

As Large Language Models (LLMs) become increasingly present in academic settings,
understanding how students evaluate and prioritise the quality of LLM-generated
explanations becomes essential. This study quantifies student judgment of LLM
explanations by examining which criteria students rely on and how they rate them
when assessing AI-generated academic content.

To address this, a student-centred evaluation framework is developed that combines
explainability criteria with multi-criterion decision making and adaptive,
dispersion-based question selection. Criterion weights are derived via the Analytic
Hierarchy Process (AHP), and Likert-scale ratings are represented through
interquartile range (IQR) intervals, preserving the range of student opinions rather
than collapsing them into a single score. An adaptive question-selection strategy
based on semantic and structural dispersion targets the most discriminative evaluation
cases, reducing annotation cost without sacrificing coverage.

The framework was validated by collecting 2{,}592 judgments expressed by nine
graduate students from a Business School in Tunisia. Applied to nine open-source LLMs
evaluated on economic reasoning tasks spanning five question categories, the
framework reveals that trustworthiness and completeness dominate student judgment,
together accounting for more than 75\% of total criterion weight. This pattern is not
captured by standard accuracy metrics. Interval-based aggregation further exposes
reasoning instability that point estimates conceal: the performance floor of the
top-ranked model (\texttt{qwen2.5:14b}) far exceeds that of the weakest model
(\texttt{codellama:7b}), and a systematic ``fluency trap'' is identified in which
models score high on understandability but low on trustworthiness.

\vspace{1em}
\noindent\textbf{Keywords:} Analytic Hierarchy Process, Human Judgment,
Inter-Rater Reliability, Interval Aggregation, Multi-Criteria Decision Making,
Large Language Models, Causal Explanation, Explainability.
